\id{МРНТИ 61.35.31}{}

\begin{header}
\swa{С.~Сатаева, В.~Бурахта, Н.~Джарасова, Т.~Өтепова, А.~Пшенбаев}{ПЕРСПЕКТИВЫ ПРИМЕНЕНИЯ ДИАТОМИТА ЗАПАДНОГО КАЗАХСТАНА В ПРОИЗВОДСТВЕ БОРОСИЛИКАТНОГО СТЕКЛА}

\tsp{1}С.~Сатаева,
\tsp{2}В.~Бурахта\envelope,
\tsp{3}Н.~Джарасова,
\tsp{1}Т.~Өтепова,
\tsp{1}А.~Пшенбаев
\end{header}

\begin{affil}
\tsp{1}Западно-Казахстанский аграрно-технический университет им. Жангир хана, Уральск, Казахстан,

\tsp{2}Западно-Казахстанский инновационно-технологический университет, Уральск, Казахстан,

\tsp{3}Уральская торгово-промышленная компания, Уральск, Казахстан

\corrauthor{Корреспондент-автор: vburakhta@mail.ru}
\end{affil}

Благодаря совокупности эксплуатационных свойств боросиликатные стекла
активно используются в изготовлении посуды для лабораторных и химических
нужд, термостойких изделий и технических компонентов, в том числе для
электронной и энергетической отраслей. В работе рассматриваются
результаты исследования химического состава местного диатомита, мела и
кварцевого песка. Проанализирована возможность использования указанных
видов минерального сырья в составе стекольной шихты при получении
боросиликатного стекла. Показано, что применение местных материалов
является перспективным направлением, позволяющим оптимизировать состав
шихты, а также снизить материальные и энергетические затраты в процессе
производства.

{\bfseries Ключевые слова}: боросиликатное стекло, кварцевый песок, мел,
диатомит, химический состав, физико-химические свойства стекла.

\begin{header}
PROSPECTS OF DIATOMITE APPLICATION IN WESTERN KAZAKHSTAN IN THE PRODUCTION OF BOROSILICATE GLASS

\tsp{1}S.~Satayeva,
\tsp{2}V.~Burakhta\envelope,
\tsp{3}N.~Dharasova,
\tsp{1}T.~Utepova,
\tsp{1}A.~Pshenbayev
\end{header}

\begin{affil}
\tsp{1}Zhangirkhan West-Kazakhstan agrarian technical university, Uralsk, Kazakhstan,

\tsp{2}West Kazakhstan Innovative Technology University, Uralsk, Kazakhstan,

\tsp{3}Ural Industrial and Trade Company, Uralsk, Kazakhstan,

e-mail: vburakhta@mail.ru
\end{affil}

Borosilicate glasses are widely used in the production of laboratory and
chemical equipment, heat-resistant and technical products, as well as in
the electronic and energy industries, due to their combination of high
heat resistance, chemical inertness, and low linear thermal expansion
coefficient. The article explores the possibility of using local
diatomite, chalk, and quartz sand as components of glass raw materials
to produce borosilicate glass. The chemical composition of the raw
materials and the physical and chemical properties of the resulting
glass are studied. It has been shown that diatomite can partially
replace quartz sand without compromising the properties of the resulting
glass. The compositions of glass charges using diatomite have been
presented. The prospects for using local mineral raw materials to
optimize the composition of the charge and reduce material and energy
costs have been established.

{\bfseries Keywords:} borosilicate glass, quartz sand, chalk, diatomite,
chemical composition, physical and che\-mical properties of glass.

\begin{header}
БОРОСИЛИКАТТЫ ШЫНЫ ӨНДІРІСІНДЕ БАТЫС ҚАЗАҚСТАН ДИАТОМИТІН ҚОЛДАНУ ПЕРСПЕКТИВАЛАРЫ

\tsp{1}С.~Сатаева,
\tsp{2}В.~Бурахта\envelope,
\tsp{3}Н.~Джарасова,
\tsp{1}Т.~Өтепова,
\tsp{1}А.~Пшенбаев
\end{header}

\begin{affil}
\tsp{1}Жәңгір хан атындағы Батыс Қазақстан аграрлық-техникалық университеті, Орал, Қазақстан,

\tsp{2}Батыс Қазақстан инновациялық-технологиялық университеті, орал, Қазақстан,

\tsp{3}«Орал сауда-өнеркәсіп компания», Орал, Қазақстан,

e-mail: vburakhta@mail.ru
\end{affil}

Боросиликатты әйнектер зертханалық және химиялық ыдыстарды, ыстыққа
төзімді және техникалық бұйымдарды өндіруде, сондай-ақ жоғары
температураға төзімділік, химиялық инерттілік және төмен сызықтық жылу
кеңею коэффициентінің үйлесімі арқасында электронды және энергетикалық
өнеркәсіпте кеңінен қолданылды. Мақалада боросиликатты шыны алу үшін
шыны шихта құрамында жергілікті диатомит, бор және кварц құмын қолдану
мүмкіндігі зерттелді. Бастапқы материалдардың химиялық құрамы және
алынған әйнектің физика-химиялық қасиеттері анықталды. Диатомит алынған
әйнектің қасиеттерін нашарлатпай, кварц құмын ішінара алмастыра
алатындығы көрсетілді. Диатомитті қолданатын шыны шихталардың құрамы
келтірілген. Шихта құрамын оңтайландыру және материалдық және
энергетикалық шығындарды азайту үшін жергілікті минералды шикізатты
пайдалану перспективасы белгіленді.

{\bfseries Түйін сөздер:} боросиликатты шыны, кварц құмы, бор, диатомит,
химиялық құрамы, шынының физика-химиялық қасиеттері.

\begin{multicols}{2}
{\bfseries Введение}. Боросиликатные стекла широко используются при
изготовлении лабораторной посуды, термостойких и технических изделий, а
также находят применение в электронной и энергетической отраслях. Такой
интерес к данным материалам обусловлен их высокой термостойкостью,
химической инертностью и пониженным коэффициентом теплового расширения
{[}1,2{]}. В качестве основного сырья для получения боросиликатных
стекол традиционно применяются кварцевые пески {[}3,4{]}. Вместе с тем
высокая температура стеклообразования и низкая реакционная способность
кварца приводят к увеличению энергетических затрат и усложняют процесс
получения однородной стекломассы {[}5-7{]}.

В этой связи актуальной является задача поиска альтернативных источников
кремнезёма, использование которых позволило бы снизить температуру
формирования стекла без ухудшения его эксплуатационных характеристик.
Результаты исследований в области стекольных технологий показывают, что
опал-кристобалитовые породы могут рассматриваться как перспективное
сырьё для стекловарения. Их отличительной особенностью является высокое
содержание аморфного диоксида кремния (до 70 \% аморфного SiO₂), а также
присутствие в составе других стеклообразующих и модифицирующих оксидов,
что способствует интенсификации процессов плавления и гомогенизации
шихты. Среди опал-кристобалитовых пород наибольший интерес для
стекольной промышленности представляют диатомиты, характеризующиеся
относительной стабильностью химического состава (табл.1). Это
обстоятельство позволяет рассматривать их как надёжный и воспроизводимый
источник кремнезёма при разработке составов стекольных шихт.
\end{multicols}

\tcap{Таблица 1 - Химический состав представителей опал-кристобалитовой группы}
\begin{longtblr}[
  label = none,
  entry = none,
]{
  width = \linewidth,
  colspec = {Q[181]Q[113]Q[100]Q[87]Q[127]Q[113]Q[113]Q[100]},
  cells = {c},
  cells = {font = \small},
  row{2} = {font=\bfseries\small},
  row{2} = {font=\bfseries\small},
  cell{1}{1} = {r=2}{font=\bfseries\small},
  cell{1}{2} = {c=7}{},
  hlines,
  vlines,
}
Название & Содержание оксидов, \% масс. &  &  &  &  &  & \\
 & SiO2 & Al2O3 & Fe2O3 & СаО & MgO & R2О & Δmпр\\
Опока & 52,1-91,4 & 2,5-15,4 & 1,0-5,0 & 0,43-17,10 & 0,08-2,48 & 0,6-4,0 & 1,5-16,8\\
Трепел & 35,3-86,7 & 2,5-11,6 & 0,3-3,4 & 0,4-31,2 & 0,2-1,6 & 0,85-2,10 & 4,5-11,6\\
Диато-мит & 77,3-83,1 & 3,3-6,0 & 1,8-3,5 & 0,3-0,6 & 0,6-1,1 & 0,8-1,5 & 3,7-8,8\\
Средняя проба & 81,98 & 5,37 & 2,67 & 0,36 & 0,80 & 1,47 & 7,90
\end{longtblr}

\begin{multicols}{2}
Значительные запасы диатомитовых пород, сосредоточенные на территории
Западного Казахстана, делают данный материал доступным и экономически
целесообразным сырьём {[}8{]}. Привлечение местных минеральных ресурсов
способствует снижению зависимости предприятий от дефицитных кварцевых
песков, уменьшению транспортных издержек и повышению устойчивости
производства. Кроме того, снижение температуры плавления и
продолжительности варки стекла при использовании активных
кремнезёмсодержащих добавок положительно сказывается на
энергопотреблении и уровне выбросов парниковых газов, что соответствует
современным требованиям устойчивого развития стекольной отрасли.

Несмотря на наличие работ, посвящённых применению диатомита в силикатных
стеклах и стеклокерамических материалах, его использование в
боросиликатных системах изучено недостаточно. В связи с этим
исследования, направленные на оценку возможности введения природного
диатомита в состав боросиликатных стекол, представляются актуальными с
точки зрения расширения сырьевой базы, повышения энергоэффективности и
развития ресурсосберегающих технологий.

Целью настоящей работы является изучение возможности получения
боросиликатного стекла с применением природного диатомита в качестве
частичного источника кремнезёма, а также оценка его влияния на основные
физико-химические свойства стекла.

В Западном регионе Казахстана, на территории Актюбинской области,
расположено месторождение Жалпак, занимающее одно из ведущих мест в мире
по запасам диатомита. Минеральный состав сырья и высокая удельная
поверхность обусловливают его повышенную реакционную способность в
процессе сплавления, что создаёт предпосылки для улучшения
технологических и эксплуатационных характеристик готовых стекольных
материалов. В связи с этим в работе исследована возможность
использования диатомита данного месторождения в качестве компонента
стекольной шихты.

{\bfseries Материалы и методы}. Элементный и массовый состав природного и
модифицированного диатомита был определён методом
рентгенофлуоресцентного анализа на спектрометре «ФОКУС-2М» (табл.2).
\end{multicols}

\tcap{Таблица 2 - Элементный и весовой составы природного и модифицированного диатомита}
\begin{longtblr}[
  label = none,
  entry = none,
]{
  cells = {c},
  cells = {font = \small},
  row{1} = {font=\bfseries\small},
  cell{1}{1} = {r=2}{},
  cell{1}{2} = {c=2}{},
  cell{1}{4} = {c=2}{},
  hlines,
  vlines,
}
Элемент & Весовой состав диатомита, \% &  & Интенсивность спектра & \\
 & природный & термически модиф. & природный & термически модиф.\\
Fe & 26,008 & 35,665 & 314,53 & 288,72\\
Ti & 1,582 & 2,933 & 11,00 & 14,93\\
K & 4,589 & 7,663 & 6,08 & 8,01\\
Ca & 0,382 & 0,548 & 1,87 & 2,17\\
Cr & 0,108 & 0,123 & 1,17 & 0,95\\
Al & 24,400 & 4,364 & 0,39 & 0,04\\
Si & 32,283 & 37,924 & 1,52 & 1,69\\
Cl & 1,081 & 1,308 & 0,38 & 0,37\\
Zn & 0,380 & 0,555 & 3,58 & 3,21\\
Cu & 0,228 & 0,061 & 1,89 & 0,31\\
As & 0,237 & 0,163 & 2,27 & 0,98\\
Rb & 0,077 & 0,348 & 0,60 & 1,70\\
Sr & 0,044 & 0,306 & 0,35 & 1,51\\
Mn & 0,055 & 0,036 & 0,72 & 0,39
\end{longtblr}

\begin{multicols}{2}
Выполнен ИК-спектральный анализ диатомита Западного Казахстана с
применением ИК Фурье-спектрометра «Varian 3100». На основании\\
литературных источников и экспериментальных исследований установлено,
что диатомит в основном представлен аморфным диоксидом кремния (SiO₂) и
карбонатом кальция (CaCO₃). Нами осуществлён обжиг диатомита в
температурном интервале 200-500 °C, в ходе которого происходит удаление
химически связанной воды из структуры минерала.

В результате данного процесса наблюдается уплотнение структуры диатомита
и существенное увеличение относительной доли диоксида кремния в образце.
ИК-спектры исходного и термически обработанного диатомита представлены
на рис.1 a, b.
\end{multicols}

\begin{figs}[Рис.1 - ИК спектры необработанного (a), прокаленного (b)
диатомита]
    \fig[0.45\textwidth]{c2/image2}[a]
    \fig[0.45\textwidth]{c2/image3}[b]
\end{figs}
\fig[0.63\textwidth]{c2/image4}[Рис.2 - Рентгенофазовый анализ кварцевого песка]

\begin{multicols}{2}
Анализ данных, приведённых на рис.1 (a, b), показывает, что полосы
поглощения в диапазоне 400--1000 см⁻¹ соответствуют асимметричным
колебаниям связей Si-O-Si; при термообработке происходит разрушение
связей Si-O-H и формирование более устойчивой кремнекислородной сетки
Si-O-Si. Диатомит, применяемый в экспериментальных исследованиях,
характеризуется высоким содержанием стеклообразующего оксида SiO₂ (до
68,88\%), а также содержит до 12,30 \% Al₂O₃ и определённое количество
модифицирующих оксидов (Fe₂O₃, R₂O, RO). Сравнительно невысокая массовая
доля SiO₂ по отношению к кварцевому песку компенсируется тем, что
основной фазой является аморфный опал, обладающий более высокой
химической активностью по сравнению с кристаллическим SiO₂.

Проведено исследование состава кварцевого песка месторождения Горбуново
Западно-Казахстанской области методом рентгено-фазового анализа {[}9{]}
(рис.2).

Результаты рентгено-фазового анализа (рис.2) свидетельствуют о том, что
образец преимущественно представлен кварцем (SiO₂), о чём говорит
высокая интенсивность основного пика, характеризующая его значительную
степень кристалличности. Пики меньшей интенсивности указывают на наличие
незначительных примесных фаз, что подтверждает высокий уровень чистоты
исследуемого материала. Результаты ранее выполненных исследований
химического состава кварцевого песка месторождения Горбуново
Западно-Казахстанской области показали, что материал характеризуется
высоким содержанием оксида кремния, при этом суммарная доля оксидов
щелочных и щелочноземельных металлов не превышает 0,3 \% {[}10,11{]}.

{\bfseries Результаты и обсуждение}. Для синтеза боросиликатного стекла
были подготовлены 3 стекольные шихты различных составов:

І состав использовался в качестве контрольного и основан на кварцевом
песке в соответствии с нормативными требованиями к химическому составу
боросиликатного стекла: \% - кварцевый песок - 42, мел - 8, борной
кислоты в пересчете на В2О3, - 23, гидрокарбоната натрия в пересчете на
Na2О - 18, Al₂O₃ - 7, MgO -- 2.

ІІ состав: \% - кварцевый песок - 42, диатомит - 10, мел - 8, борной
кислоты в пересчете на В2О3- 25, гидрокарбоната натрия в пересчете на
Na2О - 15.

ІІІ состав: \% - диатомит - 52, мел - 8, борной кислоты в пересчете на
В2О3 - 25, гидрокарбоната натрия в пересчете на Na2О - 15.

Компоненты каждой шихты тщательно смешивали и помещали в фарфоровые
тигли объёмом 100 мл. Далее тигли нагревали в муфельной печи до
температуры 1150 °С и выдерживали при этой температуре 60 минут до
получения однородного стекольного расплава.

На рис.3 представлены фотографии боросиликатного стекла, полученного из
шихты на основе местного кварцевого песка и диатомита.
\end{multicols}
\vspace{-1em}
\begin{figs}[Рис.3 - Боросиликатное стекло, полученное из шихты на основе
местного кварцевого песка и диатомита]
    \fig[0.32\textwidth]{c2/image5}[]
    \fig[0.32\textwidth]{c2/image6}[]
    \fig[0.32\textwidth]{c2/image7}[]
\end{figs}

\begin{multicols}{2}
{\bfseries Выводы}. Результаты исследований подтвердили эффективность
применения местного диатомита в составе стекольной шихты в качестве
комплексного кремнезёмсодержащего компонента при изготовлении
боросиликатного стекла. Высокое содержание аморфного диоксида кремния и
наличие функциональных оксидов позволяют диатомиту частично замещать
кварцевый песок без снижения физико-химических свойств стекла. Также
установлена возможность совместного использования местного мела и
кварцевого песка с диатомитом, что обеспечивает оптимизацию состава
шихты. Полученные данные свидетельствуют о перспективности внедрения
местного минерального сырья для повышения эффективности стекольного
производства и снижения материальных и энергетических затрат.
\end{multicols}

\begin{center}
{\bfseries Литература}
\end{center}

\begin{refs}
1. Шакиров А.А., Тлехусеж М.А. Химические процессы при изготовлении
стекла // Научное обозрение. Педагогические науки. -2019. -№ 4. -С.
93-96.

2. Макаров Г.И., Макарова Т.М. Молекулярно-динамическое моделирование
структуры боросиликатного стекла марки Е по кристаллическому
структурному шаблону // Физика и химия стекла. -2023. - T.49 (6).
-C.632-641. DOI
\href{https://doi.org/10.31857/S0132665123600498}{10.31857/S0132665123600498}.

3. Елисеева М., Гончарова С.С. Стекло: вещество или материал//Успехи в
химии и химической технологии. -2013. - № 27(5). -С.119-122.

4. Шайдуллин С.М., Никулина А.Ю., Ремизов М.Б., Козлов П.В.
Термовискозиметрические характеристики боросиликатных стекол для
перспективной области легкоплавких составов, разрабатываемых для
удаляемого малогабаритного плавителя дизайна Производственного
объединения «Маяк». // Известия Томского политехнического университета.
Инжиниринг георесурсов. -2024. -Т.335(3). - С.50-60. DOI
10.18799/24131830/2024/3/4569.

5. Шайдуллин С.М., Беланова Е.А., Козлов П.В., Ремизов М.Б.,
Дворянчикова Е.М. Отработка процесса варки боросиликатных стекол с
имитаторами компонентов ВАО и исследование их химической устойчивости //
Известия высших учебных заведений. - 2021. -Т.64 (2-2) (759). -С.
148-154. DOI 10.17223/00213411/64/2-2/148.

6. Stefanovsky S.V., Skvortsov M.V., Stefanovsky O.I., Nikonov B.S.,
Presniakov I.A., Glazkova I.S., Ptashkin A.G. Preparation and
characterization of borosilicate glass waste form for immobilization of
HLW from WWER spent nuclear fuel reprocessing. MRS Advances. -2017.
-Vol.2. - P.583-589. DOI 10.1557/adv.2016.622.

7. Бондаренко Д.О., Бондаренко Н.И., Бессмертный В.С., Изофатова Д.И.,
Дюмина П.С., Волошко Н.И. Энергосберегающая технология получения
силикат-глыбы для производства жидкого стекла // Вестник Белгородского
государственного технологического университета им. В.Г. Шухова. - 2017.
-Т.2(10) 10. - С.111-115. DOI 10.12737/article\_59cd0c60148053.59501281.

8. Смирнов П.В., Жакипбаев Б.Е., Староселец Д.А., Дерягина О.И., Баталин
Г.А., Гареев Б.И., Вергунов А.В. Диатомиты и опоки месторождений
Западного Kазахстана: литология, структурно-текстурные параметры,
потенциал использования // Известия Томского политехнического
университета. Инжиниринг георесурсов. - 2023. - №334(7). - C.187--201.
DOI:
\href{https://doi.org/10.18799/24131830/2023/7/4046}{10.18799/24131830/2023/7/4046}.

9. Анваров А.Б., Кадырова З.Р. Обогащение кварцевых песков месторождения
«Ойнакум» для синтеза высококачественного прозрачного стекла // Стекло и
керамика. -2023. -Т.96(10). -С.48-54. DOI:10.14489/glc.2023.10.
pp.048-054.

10. Бурахта В.А., Рамазанова Н.К. Применение кварцевого песка
ЗападноКазахстанской области в технологии производства тарного стекла //
Вестник ВКГТУ «Технические науки». - 2015. - №4. - С.27-30.

11. Бурахта В.А., Рамазанова Н.К. Определение состава тарного стекла
методом растрового электронного микроскопа // Материалы XII
Международной конференции «Фундаментальные и прикладные исследования в
современном мире». - 2015. - С.54-56.
\end{refs}

\begin{center}
{\bfseries References}
\end{center}

\begin{refs}
1. Shakirov A.A., Tlehusezh M.A. Himicheskie processy pri izgotovlenii
stekla // Nauchnoe obozrenie. Pedagogicheskie nauki. -2019. -№ 4. -S.
93-96. {[}in Russian{]}

2. Makarov G.I., Makarova T.M. Molekuljarno-dinamicheskoe modelirovanie
struktury borosilikatnogo stekla marki E po kristallicheskomu
strukturnomu shablonu // Fizika i himija stekla. -2023. - T.49 (6).
-C.632-641. DOI 10.31857/S0132665123600498. {[}in Russian{]}

3. Eliseeva M., Goncharova S.S. Steklo: veshhestvo ili material//Uspehi
v himii i himicheskoj tehnologii. -2013. - № 27(5). -S.119-122. {[}in
Russian{]}

4. Shajdullin S.M., Nikulina A.Ju., Remizov M.B., Kozlov P.V.
Termoviskozimetricheskie harakteristiki borosilikatnyh stekol dlja
perspektivnoj oblasti legkoplavkih sostavov, razrabatyvaemyh dlja
udaljaemogo malogabaritnogo plavitelja dizajna Proizvodstvennogo
ob\#edinenija «Majak». // Izvestija Tomskogo politehnicheskogo
universiteta. Inzhiniring georesursov. -2024. -T.335(3). - S.50-60. DOI
10.18799/24131830/2024/3/4569. {[}in Russian{]}

5. Shajdullin S.M., Belanova E.A., Kozlov P.V., Remizov M.B.,
Dvorjanchikova E.M. Otrabotka processa varki borosilikatnyh stekol s
imitatorami komponentov VAO i issledovanie ih himicheskoj ustojchivosti
// Izvestija vysshih uchebnyh zavedenij. - 2021. -T.64 (2-2) (759). -S.
148-154. DOI 10.17223/00213411/64/2-2/148. {[}in Russian{]}

6. Stefanovsky S.V., Skvortsov M.V., Stefanovsky O.I., Nikonov B.S.,
Presniakov I.A., Glazkova I.S., Ptashkin A.G. Preparation and
characterization of borosilicate glass waste form for immobilization of
HLW from WWER spent nuclear fuel reprocessing. MRS Advances. -2017.
-Vol.2. - P.583-589. DOI 10.1557/adv.2016.622.

7. Bondarenko D.O., Bondarenko N.I., Bessmertnyj V.S., Izofatova D.I.,
Djumina P.S., Voloshko N.I. Jenergosberegajushhaja tehnologija
poluchenija silikat-glyby dlja proizvodstva zhidkogo stekla // Vestnik
Belgorodskogo gosudarstvennogo tehnologicheskogo universiteta im. V.G.
Shuhova. - 2017. -T.2(10) 10. - S.111-115.
DOI 10.12737/article\_59cd0c60148053.59501281. {[}in Russian{]}

8. Smirnov P.V., Zhakipbaev B.E., Staroselec D.A., Derjagina O.I.,
Batalin G.A., Gareev B.I., Vergunov A.V. Diatomity i opoki
mestorozhdenij Zapadnogo Kazahstana: litologija, strukturno-teksturnye
parametry, potencial ispol' zovanija // Izvestija
Tomskogo politehnicheskogo universiteta. Inzhiniring georesursov. -
2023. - №334(7). - C.187--201. DOI: 10.18799/24131830/2023/7/4046.
{[}in Russian{]}

9. Anvarov A.B., Kadyrova Z.R. Obogashhenie kvarcevyh peskov
mestorozhdenija «Ojnakum» dlja sinteza vysokokachestvennogo prozrachnogo
stekla // Steklo i keramika. -2023. -T.96(10). -S.48-54.
DOI:10.14489/glc.2023.10. pp.048-054. {[}in Russian{]}

10. Burahta V.A., Ramazanova N.K. Primenenie kvarcevogo peska
ZapadnoKazahstanskoj oblasti v tehnologii proizvodstva tarnogo stekla //
Vestnik VKGTU «Tehnicheskie nauki». - 2015. - №4. - S.27-30. {[}in
Russian{]}

11. Burahta V.A., Ramazanova N.K. Opredelenie sostava tarnogo stekla
metodom rastrovogo jelektronnogo mikroskopa // Materialy XII
Mezhdunarodnoj konferencii «Fundamental' nye i prikladnye
issledovanija v sovremennom mire». - 2015. - S.54-56. {[}in Russian{]}
\end{refs}

\begin{info}
\hspace{1em}\emph{{\bfseries Сведения об авторах}}

Сатаева С.С. - кандидат химических наук, доктор PhD, ассоциированный
профессор, Западно-Казахстанский аграрно-технический университет им.
Жангир хана, Уральск, Казахстан, e-mail: sataeva\_safura@mail.ru;

Бурахта В.А. - доктор химических наук, профессор,
Западно-Казахстанский инновационно-технологический университет,
Уральск, Казахстан, e-mail: vburakhta@mail.ru;

Джарасова Н.К. - магистр технических наук, Уральская
торгово-промышленная компания, Уральск, Казахстан, e-mail:
jarassova.nurgul@mail.ru;

Өтепова Т.А. - магистрант, Западно-Казахстанский университет им.
М.Утемисова, Уральск, Казахстан, e-mail: tlekshi.utepova@mail.ru;

Пшенбаев А.Н. - магистрант, Западно-Казахстанский аграрно-технический
университет им. Жангир хана, Уральск, Казахстан, e-mail:
asan12.12.12@mail.ru.

\hspace{1em}\emph{{\bfseries Information about the authors}}

Satayeva S.S.~- Candidate of Chemical Sciences, doctor PhD, associate
professor, Zhangir Khan West Kazakhstan Agrarian and Technical
University, Uralsk, Kazakhstan, e-mail:~sataeva\_safura@mail.ru;~

Burakhta V.A.~- Doctor of Chemical Sciences, Professor, West
Kazakhstan Innovation and Technology University, Uralsk, Kazakhstan,
e-mail:~vburakhta@mail.ru;~~

Dzharasova N.K.~- Master of Engineering Sciences, Ural Industrial and
Trade Company, Uralsk, Kazakhstan, e-mail:~jarassova.nurgul@mail.ru;~

Utepova T.A. -~~Master' s student, M. Utemisov West Kazakhstan
University, Uralsk, Kazakhstan, e-mail:tlekshi.utepova@mail.ru;~~

Pshenbaev A.N. -~Master' s student, Zhangir Khan West Kazakhstan
Agrarian and Technical University, Uralsk,
Kazakhstan,~e-mail:~asan12.12.12@mail.ru.
\end{info}
